\section{Introduction}

Technical standards are fundamental to the organization of markets. By shaping compatibility, interoperability, and the size of the installed base, standards affect firms' incentives to innovate, consumers' adoption decisions, and the intensity of product-market competition. The classic literature in industrial organization has shown that network externalities may generate coordination failures, excess inertia, inefficient standardization, and ambiguous welfare effects of compatibility \citep{FarrellSaloner1985,KatzShapiro1985,KatzShapiro1986,FarrellSaloner1986,KatzShapiro1992,KatzShapiro1994}. More broadly, standards influence industry structure and economic performance through entry conditions, product variety, technological progress, and the allocation of market power \citep{DavidGreenstein1990,MatutesRegibeau1996,Shy2001,BesenFarrell1994,Tassey2017}.

These issues become even more consequential in an international setting. When standards differ across countries, firms face adaptation and compliance costs in foreign markets, and governments must decide whether to preserve national standards, recognize foreign standards, or participate in broader harmonization arrangements. A strand of the international trade literature has therefore analyzed standard recognition, mutual recognition, and strategic incompatibility as policy choices rather than purely technological facts \citep{GandalShy2001,BarrettYang2001,Costinot2008,Klimenko2009}. In particular, \citet{GandalShy2001} show that governments may have incentives to form standardization unions when network effects and conversion costs interact, while \citet{BarrettYang2001}, \citet{Costinot2008}, and \citet{Klimenko2009} emphasize that institutional design---including incompatibility, national treatment, interoperability, and mutual recognition---can systematically alter equilibrium standards and welfare. The coalition-formation aspect of this problem is also structurally related to trade-agreement models with insider-outsider tensions \citep{Riezman1985,Krishna1998}, although the object of negotiation here is standard recognition rather than tariffs.

At the same time, the economics of standards has expanded well beyond the older compatibility-versus-variety trade-off. Recent work increasingly treats standards as part of the infrastructure of innovation, regulation, and international market access. Standards and regulation affect innovation differently across market environments \citep{BlindPetersenRiillo2017}; technology standards can alter firms' product-innovation incentives \citep{FoucartLi2021}; and participation in formal standardization processes is itself shaped by strategic and institutional determinants \citep{BlindVonLaer2022,BlindKenneyLeiponenSimcoe2023}. On the trade side, empirical studies show that international standards and standard harmonization can promote cross-border exchange and market integration \citep{CloughertyGrajek2014,SchmidtSteingress2022}. More recent evidence also links international standards to export performance and innovation outcomes in a globalizing economy \citep{BlindMuench2024,WangZhu2025}. In parallel, new data efforts such as StanDat are lowering the empirical barriers to studying international standards, participation patterns, and cross-country standardization behavior \citep{Bjorkholt2026}.

A complementary institutional motivation comes from the literature on mutual
recognition agreements (MRAs), the WTO/TBT framework, and developing-country
participation in standardization \citep{WTO_TBT,ISO_UNIDO}. Those literatures
emphasize that formal recognition and market access often depend not only on
nominal membership in a common standards architecture, but also on testing
capability, certification infrastructure, regulatory implementation, and
effective participation in conformity-assessment systems. This paper abstracts
from those institutional details, but its participation-capacity parameter is
designed precisely to capture the residual compliance burden that remains when
formal recognition is granted but effective participation is incomplete. In
that sense, the model connects the compatibility literature to the
institutional question of why mutual recognition may fail to deliver fully
integrated standardization outcomes in practice.

Despite this broad and growing literature, one feature remains underexplored in formal models of international standardization. Much of the existing theory effectively treats accession to a common standard as removing the relevant cross-border friction once recognition is granted. In practice, however, formal accession and effective participation need not coincide. Countries differ not only in whether they are legally admitted to a harmonized regime, but also in whether firms can actually comply with, test, certify, and implement the standard at low cost. This distinction is particularly important when standardization is mediated by mutual recognition or bloc formation. A country may formally join a standard area yet remain a weak participant because residual compliance burdens continue to limit its contribution to the common installed base and its competitive presence in foreign markets. While the existing literature has rich insights on compatibility, standard recognition, and institutional design \citep{FarrellSaloner1985,KatzShapiro1985,GandalShy2001,BarrettYang2001,Costinot2008,Klimenko2009}, comparatively less attention has been paid to how \emph{participation-capacity asymmetry} interacts with network externalities and coalition formation in a multi-country setting.

This paper develops a simple three-country model to study that issue. The framework builds on the standards-war versus harmonization logic in the compatibility literature \citep{FarrellSaloner1985,KatzShapiro1985,KatzShapiro1986,KatzShapiro1994} and on the international standardization setting of \citet{GandalShy2001}, but introduces country-specific participation capacity. Three countries, each with one firm, choose among three regimes: a standards war ($SW$), a two-country standardization union ($SU$), and full international standardization ($IS$). Mutual recognition removes the baseline technology-gap cost inside the recognition area, but countries may still face residual compliance costs because participation capacity differs across them. In this environment, formal membership in a common standard and effective participation in that standard are distinct objects.

The model delivers three main results. First, full international standardization may maximize world welfare while failing to arise in equilibrium, because a low-capacity outsider contributes too little to the common installed base relative to the additional competitive pressure its accession creates in member-country markets. Second, stronger network effects need not monotonically favor full harmonization: over a verified parameter region with sufficiently pronounced participation-capacity asymmetry, stronger network effects can instead stabilize an exclusionary two-country standards bloc. Third, once participation capacity is endogenized, capacity-building policies and conditional transfers targeted at the low-capacity country can, on the verified compact domain used in Section~5, restore the accession margin more effectively than policies that merely preserve fragmented national standards. In that sense, the paper shifts attention from formal recognition alone to the deeper question of whether countries are able to participate effectively in a harmonized regime.

These results contribute to the literature in three ways. Relative to the classic compatibility literature, the paper introduces a tractable distinction between formal compatibility and effective participation \citep{FarrellSaloner1985,KatzShapiro1985,KatzShapiro1986,DavidGreenstein1990}. Relative to the international standardization literature, it shows how residual post-accession frictions can sustain inefficient exclusion even when mutual recognition is technologically feasible \citep{GandalShy2001,BarrettYang2001,Costinot2008,Klimenko2009}. Relative to recent empirical work on standards, trade, and innovation, it offers a mechanism through which participation asymmetries can shape the formation and persistence of standards blocs \citep{BlindPetersenRiillo2017,FoucartLi2021,BlindVonLaer2022,BlindKenneyLeiponenSimcoe2023,CloughertyGrajek2014,SchmidtSteingress2022,BlindMuench2024,WangZhu2025,Bjorkholt2026}. The resulting perspective is useful for interpreting why harmonization may fail not only because of strategic exclusion, but also because weaker participants lack the capacity to translate formal accession into effective market integration.

The policy implication is straightforward, though Section~5 makes clear that it is conditional on additional monotonicity assumptions. If fragmentation persists because low-capacity countries cannot participate effectively in a harmonized regime, then the relevant policy margin is not merely whether formal recognition is granted, but whether participation capacity can be improved. This insight links the paper's theory to contemporary discussions of standardization as part of broader innovation and regulatory infrastructure \citep{Tassey2017,BlindPetersenRiillo2017,BlindKenneyLeiponenSimcoe2023,BlindMuench2024}. In the model, targeted capacity-building and conditional transfer policies relax the actual bottleneck in the regime-selection game, whereas policies that simply preserve separate national standards do not.

The remainder of the paper is organized as follows. Section~2 presents the model. Section~3 derives firm equilibrium under $SW$, $SU$, and $IS$. Section~4 characterizes governments' regime choices under participation-capacity asymmetry. Section~5 endogenizes participation capacity through investment and studies conditional transfer policy. Section~6 concludes.
